% !TeX root = ../../Book3.tex
\section{Schwarz 引理}
\label{b3:sec:0601}

Schwarz 引理固定原点时比较欧氏模；若要比较任意两点，必须先移动圆盘坐标。归一化的代价表现为分母 \(1-\overline a z\)，它最终产生圆盘天然的非欧氏距离。

\subsection{Schwarz 引理}
\label{b3:subsec:060101}

\cref{b3:thm:02-schwarz}已经完整证明：圆盘全纯自映射固定0时，\(|f(z)|\le|z|\)。这里先看归一化为何必要。常函数 \(f(z)=a\ne0\) 在 \(z=0\) 就违反该估计；圆盘自同构 \(\phi_a\) 也不在普通欧氏距离下总为压缩。

\ProseParagraphBreak

对 \(f(0)=0\) 的函数，商 \(f(z)/z\) 在0补为导数后全纯且模不超过1。这个商可以再次使用最大模或 Schwarz 引理，给高阶零点的估计。若零点阶数至少为 \(m\)，逐次除以 \(z\) 得 \(|f(z)|\le|z|^m\)。每次使用时，商要么为单位常数，要么仍把盘映入开盘；若 \(m\ge2\)，单位常数与剩余零点条件矛盾，因此迭代合法。

\ProseParagraphBreak

例如 \(f(z)=z^m\) 达到这一上界。条件只是零点阶数，不要求函数在整个盘单叶；\(m>1\) 的幂映射反而提供最简单的非单叶例子。

Schwarz 引理的证明把一个零点从函数中除掉：\(h(z)=f(z)/z\) 在0可去。若边界圆 \(|z|=r<1\) 上只有 \(|f|<1\)，便先得到 \(|h|\le1/r\)，再令 \(r\uparrow1\)。这个顺序不使用单位圆上的边界值。因而对没有边界连续延拓的圆盘函数同样成立。把“圆周上 \(|f|\le1\)”写在单位圆周而不检查函数是否有定义，会无故增加一个并不存在的假设。

\ProseParagraphBreak
取等时除零点后的函数在内点达到最大模，所以必为单位常数。由此，一个看似只是数量估计的结论同时识别了全部极端映射。许多复分析极值论证都是这一模式：先把约束转换成一个有界全纯函数，再将等号交给最大模原理。等号条件因此不能只靠形式求导获得。

\subsection{导数形式}
\label{b3:subsec:060102}

导数形式 \(|f'(0)|\le1\) 衡量原点的局部伸缩。若定义域是 \(D(a,R)\)，目标是 \(D(b,S)\)，且 \(f(a)=b\)，先置 \(F(\zeta)=(f(a+R\zeta)-b)/S\)，便得到
\[
 |f(z)-b|\le\frac SR|z-a|,\qquad |f'(a)|\le\frac SR.
\]
这两个尺度因子分别来自源坐标与目标坐标，不能只改变盘半径而保持原来的常数1。

\ProseParagraphBreak

若不固定目标中心，可以先接目标盘自同构，便获得稍后要证明的更精确导数估计。在圆盤中心 \(a=0\) 时，它变成 \(|f'(0)|\le1-|f(0)|^2\)，比直接 Cauchy 估计 \(|f'(0)|\le1\) 严格反映目标点靠近边界时允许的伸缩减少。

\begin{proposition}\label{b3:prop:06-blaschke-division}
设 \(f:\mathbb D\rightarrow\mathbb D\) 全纯，\(f(a)=0\)。则 \(f(z)/\phi_a(z)\) 在 \(a\) 可去，且模不超过1。若 \(a_1,\ldots,a_m\) 是按重数列出的有限组零点，则
\[
 |f(z)|\le\prod_{j=1}^m|\phi_{a_j}(z)|.
\]
\end{proposition}
\begin{proof}
将 \(a\) 送至0后用 Schwarz 引理，再撤销自同构，得到第一式。商全纯且模不超过1；若商不是单位常数，开映射定理使其值仍在圆盘内，可以继续对下一个零点除法。若成为单位常数，则不再有剩余零点。逐次相乘得到结论。
\end{proof}
有限乘积 \(B(z)=\eta\prod_j\phi_{a_j}(z)\)、\(|\eta|=1\)，称为有限 Blaschke 乘积。它把圆盘映入圆盘，单位圆上的模为1，在盘内恰有规定的零点。上述估计说明在给定这些零点的圆盘函数中，\(|B(z)|\) 是逐点可达到的上界。若某个不属于零点的内点取等，最后的商由最大模原理为单位常数，故 \(f\) 必是这样的有限乘积。

\ProseParagraphBreak
例如给定零点0和 \(a\)，便有 \(|f(z)|\le |z|\,|\phi_a(z)|\)，比只用0的 Schwarz 估计严格。若两个零点重合于0，就恢复 \(|f(z)|\le|z|^2\)。零点的信息能累积，但无限零点必须另外讨论无穷乘积的收敛，本章不以一个未经验证的无穷乘积代替有限除法。

\subsection{等号情形}
\label{b3:subsec:060103}

Schwarz 引理的等号结论来自最大模原理，不是只由导数数值猜出的。若 \(f(0)=0\) 且 \(|f'(0)|=1\)，商 \(f(z)/z\) 的模在内部达到1，故它是单位常数。若某个非零点 \(z_0\) 满足 \(|f(z_0)|=|z_0|\)，同一商在 \(z_0\) 达到最大模，也得到旋转。

\ProseParagraphBreak

只在趋近边界的点列上渐近取等，不会迫使旋转。比如 \(f(z)=z^2\)，有 \(|f(r)|/r=r\rightarrow1\)，但 \(f'(0)=0\)。刚性要求一个内部点上的精确等号。定量的接近等号问题需要另外估计，不能把边界极限直接代入内部等号情形。

导数的尺度也可以直接读出。若 \(f:D(a,R)\rightarrow D(b,S)\) 且 \(f(a)=b\)，则 \(|f'(a)|\le S/R\)。当取等时 \(f(z)=b+(S/R)\eta(z-a)\)。这表明全纯映射并不普遍压缩欧氏长度：目标圆盘的半径比源圆盘大时，最简单的线性映射已经放大长度。Schwarz 引理压缩的是归一化坐标中的长度。

\ProseParagraphBreak
若 \(|f(z)|\le M\) 但没有指定中心值为0，只能由 Cauchy 估计得到 \(|f'(a)|\le M/R\)。一旦知道 \(f(a)=c\)，Schwarz--Pick 会改进这一上界为 \((M^2-|c|^2)/(MR)\)。函数值接近目标边界时可用的导数空间变小；在极限 \(|c|=M\) 时最大模原理使函数常值。下面的分式归一化正是把这个现象定量化。

\subsection{圆盘自同构}
\label{b3:subsec:060104}

圆盘自同构已在\cref{b3:prop:05-disk-automorphisms}分类。记 \(\phi_a(z)=(z-a)/(1-\overline a z)\)，则
\[
 \phi_a^{-1}=\phi_{-a},\qquad
 \phi_a'(z)=\frac{1-|a|^2}{(1-\overline a z)^2},\qquad
 1-|\phi_a(z)|^2=\frac{(1-|a|^2)(1-|z|^2)}{|1-\overline a z|^2}.
\]
前两式可直接解方程与求导，第三式由两边平方模通分得出。三式共同给出将任意点送到0后的导数换算。

\ProseParagraphBreak

圆盘自同构未必固定任何内部点。例如 \(\phi_a\) 的固定点方程在 \(a\ne0\) 时为 \(\overline a z^2=a\)，两个解都在单位圆上。因而不能把所有自同构都用一个内部固定点旋转表示；旋转分类只适用于已经归一化固定0的自同构。

圆盘自同构的导数恒等式为
\[
 |T'(z)|=\frac{1-|T(z)|^2}{1-|z|^2}.
\]
可对 \(\phi_a\) 直接求导，再利用旋转与复合链式法则验证。它说明同一自同构可能在某处放大欧氏长度、在另一处缩小长度，却处处保持经边界距离校正的尺度。例如 \(T(z)=(z+a)/(1+az)\)、\(0<a<1\)，在0的导数为 \(1-a^2<1\)，在 \(-a\) 的导数为 \(1/(1-a^2)>1\)。两处都在单位盘内部，因而欧氏导数小于1不能作为圆盘映射的全局性质。

\ProseParagraphBreak
这条恒等式同时保证不同归一化坐标给出的 Schwarz--Pick 结论一致。若在源域或目标域更换圆盘坐标，各导数多出的自同构因子恰由 \(1-|z|^2\) 的变化抵消。后面的度量构造不是另加一个任意分母，而是把这里已发现的坐标不变量写成长度形式。
